%-------------------------------------------------------------------------------
% paper/paper.tex — the LaTeX twin of paper.md (the JOSS submission artifact).
% JOSS itself ONLY accepts paper.md; this .tex produces the same article for
% preprint / co-author circulation, sharing the same paper.bib.
% Build:  pdflatex paper && bibtex paper && pdflatex paper && pdflatex paper
%-------------------------------------------------------------------------------
\documentclass[11pt,a4paper]{article}

\usepackage[T1]{fontenc}
\usepackage[utf8]{inputenc}
\usepackage{lmodern}
\usepackage[english]{babel}
\usepackage[margin=2.6cm]{geometry}
\usepackage{microtype}
\usepackage[numbers,sort&compress]{natbib}
\usepackage{xcolor}
\definecolor{linkink}{RGB}{20,60,110}
\usepackage[colorlinks=true, linkcolor=linkink,
            citecolor=linkink, urlcolor=linkink]{hyperref}

\title{Choupo: a glass-box chemical process simulator for teaching
       transport phenomena and process engineering}
% TODO(Vítor): decidir a lista de autores (a autoria do paper é decisão tua,
% não deriva do copyright do código) e acrescentar o ORCID.
\author{Vítor Geraldes\\[2pt]
  \small Department of Chemical Engineering, Instituto Superior Técnico,\\
  \small Universidade de Lisboa, Portugal}
\date{20 July 2026}

\begin{document}
\maketitle

\begin{abstract}
Choupo is an open-source, educational chemical process simulator written in
C++17 with zero external dependencies, in which every numerical method is
readable, every model choice is announced at run time, and conservation
laws are first-class teaching artefacts.  Cases are plain-text
dictionaries; the graphical interface is a runner and visualiser built on
the same solvers compiled to WebAssembly.  The current release ships four
solver binaries, 49 unit-operation models, a curated thermodynamic
catalogue and 288 regression-validated tutorial cases.
\end{abstract}

\section*{Summary}

Choupo is an open-source, educational chemical process simulator written in
C++17 with \textbf{zero external dependencies}: every numerical method it
uses --- Newton--Raphson in one and many dimensions, Gauss elimination,
Levenberg--Marquardt regression, Runge--Kutta and Rosenbrock integration,
Wegstein acceleration, Nelder--Mead search and Michelsen's tangent-plane
stability test \citep{michelsen1982} --- is implemented in readable source
the student can open, step through, and modify.  Cases are plain-text
dictionaries on disk, in a file-first layout in which the flowsheet file
declares only topology while every stream's state lives in its own file;
the graphical interface is a runner and visualiser (the same solvers
compiled to WebAssembly run in the browser), never a black-box editor.

The current release ships four solver binaries by problem class
(steady-state flowsheeting, batch campaigns driven by time-triggered
recipes, dynamic simulation with feedback control, and a property
workbench), 49 unit-operation models, a curated thermodynamic catalogue
(247 components, 205 Henry's-law pairs, Pitzer and electrolyte-NRTL
parameter sets), and 288 runnable tutorial cases validated by a regression
suite with golden-master checks on every release.

\section*{Statement of need}

Commercial process simulators dominate both industry and the classroom, but
they are pedagogically opaque: the student specifies a column and receives
converged profiles without ever seeing a residual, a Jacobian, or the
reference state of an enthalpy.  Existing open-source alternatives
prioritise breadth and industrial workflows over readability of the
computation itself.  Choupo takes the opposite stance --- \emph{transparency
over breadth}: at the default verbosity every Newton iteration, every
K-value and every convergence aid is printed and attributed, and the engine
refuses to silently substitute defaults for missing data (a missing binary
pair is announced; a missing formation enthalpy makes an energy balance
refuse rather than fabricate a zero).

Three design decisions serve that goal.  First, \textbf{conservation is the
curriculum}: every converged steady run emits a plant-boundary element
balance (atoms of each element in versus out --- the true invariant of a
reacting system), batch campaigns carry material and energy ledgers whose
closures are exact state differences on the elements-of-formation datum,
and dynamic runs integrate a conservation ledger on accepted integrator
steps.  Claims are tri-state (full, partial, unavailable) with named
reasons, never silent zeros.  Second, \textbf{thermodynamics is declared,
not defaulted}: a case carries a single thermophysical-system dictionary
that names its equilibrium formulation (activity--fugacity, dilute-solution
Henry, a single cubic equation of state serving both phases
\citep{soave1972,peng1976}, or aqueous electrolyte), and the run log
announces every assembled model and parameter source.  The model library
spans NRTL \citep{renon1968}, UNIQUAC \citep{abrams1975}, predictive UNIFAC
\citep{fredenslund1975} and COSMO-SAC \citep{lin2002,bell2019}, Pitzer
\citep{pitzer1973} and electrolyte-NRTL \citep{chen1986} electrolytes, and
a validated non-associating PC-SAFT core \citep{gross2001}.  Third,
\textbf{cases are self-contained}: an importer seals each case's property
records into the case directory under a checksummed manifest, so a tutorial
archived with a thesis reruns byte-for-byte years later with no
installation catalogue present.

Choupo is an independent personal project, conceived, self-funded and
developed by the author in his free time, on personal equipment; it is not
an institutional product.  It is designed for teaching transport phenomena
and process engineering --- and for the LLM era of case authoring: the
documentation ships a machine-oriented reference so a student can author
dictionaries with an assistant while the engine remains the sole arbiter
of the physics.

\section*{Acknowledgements}

Substantial parts of the initial implementation, documentation and tutorial
corpus were produced with assistance from Anthropic Claude; the published
project is human-curated, reviewed and maintained by the authors.  The
\emph{Choupo} name is a trademark of TalentGround Lda.; the code is
GPL-3.0-or-later.

\bibliographystyle{unsrtnat}
\bibliography{paper}

\end{document}
